\subsection{Groups, Subgroups, and Homomorphisms}

Groups are the algebraic language of reversible operations. They appear here as the additive structure underlying rings and modules, and later as symmetry groups acting on mathematical objects.

\begin{definition} \label{def:group-and-abelian-group} \hlt{Groups}\\
A \hlt{group} is a nonempty set $G$ with a binary operation $(a,b)\mapsto ab$ such that:
\begin{enumerate}[label=(\roman*)]
\setlength{\itemsep}{0pt}
\item $(ab)c=a(bc)$ for all $a,b,c\in G$;
\item there is an element $e\in G$ with $eg=ge=g$ for all $g\in G$;
\item for each $g\in G$, there is $g^{-1}\in G$ with $gg^{-1}=g^{-1}g=e$.
\end{enumerate}
The group is \hlt{abelian} if $ab=ba$ for all $a,b\in G$.\\
In an abelian group the operation is often written additively, with identity $0$ and inverse $-a$.
\end{definition}

\begin{example} \label{ex:basic-groups-from-linear-algebra} \hlt{Basic Groups}
\begin{enumerate}[label=(\roman*)]
\setlength{\itemsep}{0pt}
\item The bijections $S\to S$ form a group under composition; it is usually nonabelian.
\item $\operatorname{Mat}(m,n,\mathbb{F})$ is an abelian group under matrix addition.
\item $\operatorname{GL}_n(\mathbb{F})$, the invertible $n\times n$ matrices, is a group under multiplication.
\item $\Z/n\Z$ is a finite abelian group under addition modulo $n$.
\end{enumerate}
\end{example}

\begin{definition} \label{def:subgroup} \hlt{Subgroup}\\
A \hlt{subgroup} of a group $G$ is a nonempty subset $H\subseteq G$ that is itself a group under operation inherited from $G$.
\end{definition}

\begin{proposition} \label{prop:subgroup-test} \hlt{Subgroup Test}\\
Let $G$ be a group and let $H\subseteq G$ be nonempty. Then $H\leq G$ $\iff$ $ab^{-1}\in H \ \forall a,b\in H$.
\begin{proof}
\begin{enumerate}[label=\roman*.]
\setlength{\itemsep}{0pt}
\item[$(\Rightarrow)$] If $H\leq G$, then $b^{-1}\in H$ and $ab^{-1}\in H$ for all $a,b\in H$.
\item[$(\Leftarrow)$] Assume $ab^{-1}\in H$ for all $a,b\in H$. Choose $h\in H$. Then $e_G=hh^{-1}\in H$. For $b\in H$, $b^{-1}=e_Gb^{-1}\in H$.\\
For $a,c\in H$, since $c^{-1}\in H$, we have $ac=a(c^{-1})^{-1}\in H$. Hence $H\leq G$.
\end{enumerate}
\hlt{Structural Commentary}: Closure under one expression packages identity, inverses, and products.\\
\hlt{Key Insight}: In a group, subtraction-style closure is enough.
\end{proof}
\end{proposition}

\begin{definition} \label{def:cyclic-group-and-generated-subgroup} \label{def:group-homomorphism-kernel-image} \hlt{Generated Groups and Homomorphisms}\\
Let $G,H$ be groups.
\begin{enumerate}[label=(\roman*)]
\setlength{\itemsep}{0pt}
\item For $a\in G$, the \hlt{subgroup generated by $a$} is $\langle a\rangle\defeq\{a^k:k\in\Z\}\leq G$.
\item $G$ is \hlt{cyclic} if $G=\langle a\rangle$ for some $a\in G$.\\
If it exists, $\operatorname{ord}(a)\defeq\min\{n\in\N_{>0}:a^n=e_G\}$; then $\langle a\rangle=\{e_G,a,\ldots,a^{\operatorname{ord}(a)-1}\}$.
\item A \hlt{homomorphism} $\varphi:G\to H$ satisfies $\varphi(ab)=\varphi(a)\varphi(b)$ for all $a,b\in G$.
\item For a homomorphism $\varphi:G\to H$, $\ker\varphi\defeq\{g\in G:\varphi(g)=e_H\}$; and $\operatorname{im}\varphi\defeq\varphi(G)$.\\
An \hlt{isomorphism} is a bijective homomorphism.
\end{enumerate}
\end{definition}

\begin{proposition} \label{prop:group-homomorphism-basic-properties} \hlt{Basic Homomorphism Properties}\\
Let $\varphi:G\to H$ be a group homomorphism. Then:
\begin{enumerate}[label=(\roman*)]
\setlength{\itemsep}{0pt}
\item $\varphi(e_G)=e_H$;
\item $\varphi(g^{-1})=\varphi(g)^{-1}$ for all $g\in G$;
\item $\ker\varphi\leq G$;
\item $\operatorname{im}\varphi\leq H$.
\end{enumerate}
\begin{proof}
\begin{enumerate}[label=(\roman*)]
\setlength{\itemsep}{0pt}
\item Since $\varphi(e_G)=\varphi(e_Ge_G)=\varphi(e_G)^2$, cancellation gives $\varphi(e_G)=e_H$.
\item For $g\in G$, $\varphi(g)\varphi(g^{-1})=\varphi(gg^{-1})=e_H$, hence $\varphi(g^{-1})=\varphi(g)^{-1}$.
\item If $a,b\in\ker\varphi$, then $\varphi(ab^{-1})=\varphi(a)\varphi(b)^{-1}=e_H$, so $\ker\varphi\leq G$ by Proposition \ref{prop:subgroup-test}.
\item If $x=\varphi(a)$ and $y=\varphi(b)$ lie in $\operatorname{im}\varphi$, then $xy^{-1}=\varphi(ab^{-1})\in\operatorname{im}\varphi$, so $\operatorname{im}\varphi\leq H$ by Proposition \ref{prop:subgroup-test}.
\end{enumerate}
\hlt{Structural Commentary}: Homomorphisms preserve all group structure once they preserve multiplication.\\
\hlt{Key Insight}: $\varphi(ab)=\varphi(a)\varphi(b)$ preserves $e$, inverses, kernels, and images.
\end{proof}
\end{proposition}

\begin{proposition} \label{prop:group-homomorphism-injective-kernel} \hlt{Kernel Criterion for Injectivity}\\
Let $\varphi:G\to H$ be a group homomorphism. Then $\varphi$ injective $\iff$ $\ker\varphi=\{e_G\}$.
\begin{proof}
\begin{enumerate}[label=\roman*.]
\setlength{\itemsep}{0pt}
\item[$(\Rightarrow)$] If $\varphi$ injective, $g\in\ker\varphi$, then $\varphi(g)=e_H=\varphi(e_G)$ by Proposition \ref{prop:group-homomorphism-basic-properties}, so $g=e_G$. Hence $\ker\varphi=\{e_G\}$.
\item[$(\Leftarrow)$] If $\ker\varphi=\{e_G\}$ and $\varphi(a)=\varphi(b)$, then $\varphi(ab^{-1})=\varphi(a)\varphi(b)^{-1}=e_H$, so $ab^{-1}=e_G$ and $a=b$.
\end{enumerate}
\hlt{Structural Commentary}: Injectivity is detected by the identity fiber $\ker\varphi$.\\
\hlt{Key Insight}: The kernel measures exactly the failure of injectivity.
\end{proof}
\end{proposition}
